\chapter{DF-060 Audit viewport masking and scrollable inspection panes family across real surfaces}\label{record-100}

\noindent{\footnotesize\color{BookMuted}\path{docs/design/DF-060-audit-viewport-masking-and-scrollable-inspection-panes-family-across-real-surfaces.md}\par}\smallskip

\begin{RecordMeta}
title: DF-060 Audit viewport masking and scrollable inspection panes family across real surfaces
legend: DF
lane: design
priority: medium
keywords:
  - component-audit
  - dogfood
  - lowerings
  - viewport

\end{RecordMeta}

\section{Framing}\label{record-100-framing}

Viewport masking is the family that proves overflow is owned by bounded panes, not by callers slicing strings after rendering. \texttt{v\allowbreak{}i\allowbreak{}e\allowbreak{}w\allowbreak{}p\allowbreak{}o\allowbreak{}r\allowbreak{}t\allowbreak{}S\allowbreak{}u\allowbreak{}r\allowbreak{}f\allowbreak{}a\allowbreak{}c\allowbreak{}e\allowbreak{}(\allowbreak{})} is the base masking primitive for structured surfaces, \texttt{p\allowbreak{}a\allowbreak{}g\allowbreak{}e\allowbreak{}r\allowbreak{}S\allowbreak{}u\allowbreak{}r\allowbreak{}f\allowbreak{}a\allowbreak{}c\allowbreak{}e\allowbreak{}(\allowbreak{})} adds long-document position context, and \texttt{f\allowbreak{}o\allowbreak{}c\allowbreak{}u\allowbreak{}s\allowbreak{}A\allowbreak{}r\allowbreak{}e\allowbreak{}a\allowbreak{}S\allowbreak{}u\allowbreak{}r\allowbreak{}f\allowbreak{}a\allowbreak{}c\allowbreak{}e\allowbreak{}(\allowbreak{})} adds workspace focus ownership. Their string counterparts remain explicit text-lowering paths.

This audit verifies the DOGFOOD \texttt{v\allowbreak{}i\allowbreak{}e\allowbreak{}w\allowbreak{}p\allowbreak{}o\allowbreak{}r\allowbreak{}t\allowbreak{}-\allowbreak{}s\allowbreak{}u\allowbreak{}r\allowbreak{}f\allowbreak{}a\allowbreak{}c\allowbreak{}e} story against real story rendering, component-family doctrine, and mode-aware lowerings. It also closes the story drift where DOGFOOD taught only the base viewport mask even though the family doc already included pager and focus-area panes.

\section{Sponsored Users}\label{record-100-sponsored-users}

\begin{itemize}
\tightlist
\item
  TUI app authors wrapping existing surfaces in bounded scroll regions.
\item
  Docs readers checking the distinction between masking, paging, and focused inspection panes.
\item
  Maintainers who need one focused regression tying viewport story metadata, variants, and design-system guidance to real output.
\end{itemize}

\section{Hills}\label{record-100-hills}

\begin{enumerate}
\def\labelenumi{\arabic{enumi}.}
\tightlist
\item
  A builder can inspect DOGFOOD and see viewport mask, structured stack, pager, and focused pane variants in one family.
\item
  A reader can switch every variant through interactive, static, pipe, and accessible profiles while preserving scroll and focus context.
\item
  A maintainer can run one cycle test that proves viewport-family story metadata and family guidance have not drifted away from runtime behavior.
\end{enumerate}

\section{Playback Questions}\label{record-100-playback-questions}

\begin{itemize}
\tightlist
\item
  Does DOGFOOD expose the family through the \texttt{v\allowbreak{}i\allowbreak{}e\allowbreak{}w\allowbreak{}p\allowbreak{}o\allowbreak{}r\allowbreak{}t\allowbreak{}-\allowbreak{}m\allowbreak{}a\allowbreak{}s\allowbreak{}k\allowbreak{}i\allowbreak{}n\allowbreak{}g\allowbreak{}-\allowbreak{}a\allowbreak{}n\allowbreak{}d\allowbreak{}-\allowbreak{}s\allowbreak{}c\allowbreak{}r\allowbreak{}o\allowbreak{}l\allowbreak{}l\allowbreak{}a\allowbreak{}b\allowbreak{}l\allowbreak{}e\allowbreak{}-\allowbreak{}i\allowbreak{}n\allowbreak{}s\allowbreak{}p\allowbreak{}e\allowbreak{}c\allowbreak{}t\allowbreak{}i\allowbreak{}o\allowbreak{}n\allowbreak{}-\allowbreak{}p\allowbreak{}a\allowbreak{}n\allowbreak{}e\allowbreak{}s} coverage id?
\item
  Do variants cover pure viewport masking, structured masking, pager status, and focus-area chrome?
\item
  Do rich/static previews keep bounded pane chrome visible?
\item
  Do pipe lowerings preserve sequential content and scroll/focus context without pretending hidden regions are interactive?
\item
  Do accessible lowerings linearize pane content with explicit position or focus facts?
\item
  Do component-family docs still cover viewport, pager, focus area, and the row-aware holdout boundary?
\end{itemize}

\section{Requirements}\label{record-100-requirements}

\begin{itemize}
\tightlist
\item
  Keep the existing \texttt{v\allowbreak{}i\allowbreak{}e\allowbreak{}w\allowbreak{}p\allowbreak{}o\allowbreak{}r\allowbreak{}t\allowbreak{}-\allowbreak{}s\allowbreak{}u\allowbreak{}r\allowbreak{}f\allowbreak{}a\allowbreak{}c\allowbreak{}e} story identity and coverage family id.
\item
  Render every viewport-family story variant in every canonical story profile.
\item
  Treat scrollbars, gutters, and box drawing as visual-profile chrome only.
\item
  Keep pipe and accessible lowerings scroll-context-first and semantic-preserving.
\item
  Confirm \texttt{d\allowbreak{}o\allowbreak{}c\allowbreak{}s\allowbreak{}/\allowbreak{}d\allowbreak{}e\allowbreak{}s\allowbreak{}i\allowbreak{}g\allowbreak{}n\allowbreak{}-\allowbreak{}s\allowbreak{}y\allowbreak{}s\allowbreak{}t\allowbreak{}e\allowbreak{}m\allowbreak{}/\allowbreak{}c\allowbreak{}o\allowbreak{}m\allowbreak{}p\allowbreak{}o\allowbreak{}n\allowbreak{}e\allowbreak{}n\allowbreak{}t\allowbreak{}-\allowbreak{}f\allowbreak{}a\allowbreak{}m\allowbreak{}i\allowbreak{}l\allowbreak{}i\allowbreak{}e\allowbreak{}s\allowbreak{}.\allowbreak{}m\allowbreak{}d} still matches the runtime story posture for viewport, pager, and focus-area primitives.
\end{itemize}

\section{Acceptance Criteria}\label{record-100-acceptance-criteria}

\begin{itemize}
\tightlist
\item
  \texttt{t\allowbreak{}e\allowbreak{}s\allowbreak{}t\allowbreak{}s\allowbreak{}/\allowbreak{}c\allowbreak{}y\allowbreak{}c\allowbreak{}l\allowbreak{}e\allowbreak{}s\allowbreak{}/\allowbreak{}D\allowbreak{}F\allowbreak{}-\allowbreak{}0\allowbreak{}6\allowbreak{}0\allowbreak{}/\allowbreak{}v\allowbreak{}i\allowbreak{}e\allowbreak{}w\allowbreak{}p\allowbreak{}o\allowbreak{}r\allowbreak{}t\allowbreak{}-\allowbreak{}m\allowbreak{}a\allowbreak{}s\allowbreak{}k\allowbreak{}i\allowbreak{}n\allowbreak{}g\allowbreak{}-\allowbreak{}f\allowbreak{}a\allowbreak{}m\allowbreak{}i\allowbreak{}l\allowbreak{}y\allowbreak{}-\allowbreak{}a\allowbreak{}u\allowbreak{}d\allowbreak{}i\allowbreak{}t\allowbreak{}.\allowbreak{}t\allowbreak{}e\allowbreak{}s\allowbreak{}t\allowbreak{}.\allowbreak{}t\allowbreak{}s} proves the cycle doc carries the modern playback sections.
\item
  The test locates the DOGFOOD \texttt{v\allowbreak{}i\allowbreak{}e\allowbreak{}w\allowbreak{}p\allowbreak{}o\allowbreak{}r\allowbreak{}t\allowbreak{}S\allowbreak{}u\allowbreak{}r\allowbreak{}f\allowbreak{}a\allowbreak{}c\allowbreak{}e\allowbreak{}(\allowbreak{})\allowbreak{} \allowbreak{}/\allowbreak{} \allowbreak{}p\allowbreak{}a\allowbreak{}g\allowbreak{}e\allowbreak{}r\allowbreak{}S\allowbreak{}u\allowbreak{}r\allowbreak{}f\allowbreak{}a\allowbreak{}c\allowbreak{}e\allowbreak{}(\allowbreak{})\allowbreak{} \allowbreak{}/\allowbreak{} \allowbreak{}f\allowbreak{}o\allowbreak{}c\allowbreak{}u\allowbreak{}s\allowbreak{}A\allowbreak{}r\allowbreak{}e\allowbreak{}a\allowbreak{}S\allowbreak{}u\allowbreak{}r\allowbreak{}f\allowbreak{}a\allowbreak{}c\allowbreak{}e\allowbreak{}(\allowbreak{})} story and verifies expected variants.
\item
  Every variant renders non-empty output in every documented profile.
\item
  Rich/static output includes structured visual containment; pipe and accessible output do not depend on box drawing.
\item
  Pipe and accessible output preserve viewport scroll, pager line, and focus pane context.
\item
  Component-family docs retain viewport, pager, focus area, content guidance, lowering guidance, and the domain-specific row-aware boundary.
\end{itemize}

\section{Implementation Outline}\label{record-100-implementation-outline}

\begin{enumerate}
\def\labelenumi{\arabic{enumi}.}
\tightlist
\item
  Extend the existing \texttt{v\allowbreak{}i\allowbreak{}e\allowbreak{}w\allowbreak{}p\allowbreak{}o\allowbreak{}r\allowbreak{}t\allowbreak{}-\allowbreak{}s\allowbreak{}u\allowbreak{}r\allowbreak{}f\allowbreak{}a\allowbreak{}c\allowbreak{}e} DOGFOOD story with pager and focused pane variants.
\item
  Add mode-aware constrained lowerings for the story so pipe/accessibility profiles expose scroll/focus facts as text rather than visual chrome.
\item
  Add a DF-060 cycle test that renders every variant through the shared story protocol.
\item
  Move the backlog note into \texttt{d\allowbreak{}o\allowbreak{}c\allowbreak{}s\allowbreak{}/\allowbreak{}d\allowbreak{}e\allowbreak{}s\allowbreak{}i\allowbreak{}g\allowbreak{}n\allowbreak{}/} and update the v6 lane pointer.
\end{enumerate}

\section{Drift Check}\label{record-100-drift-check}

Runtime support already existed for the full family. The drift was in DOGFOOD: the story represented only \texttt{v\allowbreak{}i\allowbreak{}e\allowbreak{}w\allowbreak{}p\allowbreak{}o\allowbreak{}r\allowbreak{}t\allowbreak{}S\allowbreak{}u\allowbreak{}r\allowbreak{}f\allowbreak{}a\allowbreak{}c\allowbreak{}e\allowbreak{}(\allowbreak{})} while the design-system family also documented \texttt{p\allowbreak{}a\allowbreak{}g\allowbreak{}e\allowbreak{}r\allowbreak{}S\allowbreak{}u\allowbreak{}r\allowbreak{}f\allowbreak{}a\allowbreak{}c\allowbreak{}e\allowbreak{}(\allowbreak{})} and \texttt{f\allowbreak{}o\allowbreak{}c\allowbreak{}u\allowbreak{}s\allowbreak{}A\allowbreak{}r\allowbreak{}e\allowbreak{}a\allowbreak{}S\allowbreak{}u\allowbreak{}r\allowbreak{}f\allowbreak{}a\allowbreak{}c\allowbreak{}e\allowbreak{}(\allowbreak{})}.

This slice fixes that by extending the story instead of filing follow-on debt. No runtime API changes were required.

\section{Playback}\label{record-100-playback}

\begin{itemize}
\tightlist
\item
  RED: the release lane had only a backlog note, and the DOGFOOD story did not visibly exercise pager or focused-pane variants.
\item
  GREEN: the new cycle test renders \texttt{d\allowbreak{}o\allowbreak{}c\allowbreak{}u\allowbreak{}m\allowbreak{}e\allowbreak{}n\allowbreak{}t}, \texttt{s\allowbreak{}t\allowbreak{}r\allowbreak{}u\allowbreak{}c\allowbreak{}t\allowbreak{}u\allowbreak{}r\allowbreak{}e\allowbreak{}d\allowbreak{}-\allowbreak{}s\allowbreak{}t\allowbreak{}a\allowbreak{}c\allowbreak{}k}, \texttt{p\allowbreak{}a\allowbreak{}g\allowbreak{}e\allowbreak{}r\allowbreak{}-\allowbreak{}w\allowbreak{}i\allowbreak{}n\allowbreak{}d\allowbreak{}o\allowbreak{}w}, and \texttt{f\allowbreak{}o\allowbreak{}c\allowbreak{}u\allowbreak{}s\allowbreak{}e\allowbreak{}d\allowbreak{}-\allowbreak{}p\allowbreak{}a\allowbreak{}n\allowbreak{}e} variants across interactive, static, pipe, and accessible profiles.
\item
  Rich/static previews preserve scrollbars, status lines, focus gutters, and nested surface masking.
\item
  Pipe previews preserve viewport scroll, pager line, and focus context without box drawing.
\item
  Accessible previews linearize the same pane facts in reading order.
\end{itemize}

\section{Retrospective}\label{record-100-retrospective}

The runtime did not need another overflow abstraction. The useful release work was making DOGFOOD tell the whole family story: mask structured surfaces, page long linear documents, and mark focused inspection panes without letting scroll chrome become semantic content.

